\section{Methodology}
\label{sec:methodology}

\subsection{Conceptual Framework: Two Axes, Four Variants}

Every experiment in this thesis compares encodings of the same problem that differ only in how the domain-specific heuristic is represented and interpreted. The comparison is structured by two orthogonal axes.

The \emph{approach} axis distinguishes where the heuristic is evaluated. In the ground-and-solve approach, the heuristic is written with native \mintinline{text}|#heuristic| directives; gringo expands them before search, and clasp's Domain heuristic consumes the resulting ground objects. In the lazy approach, the heuristic remains a handful of ground facts; a custom propagator reads them at initialization and evaluates them at every decision point against the partial assignment.

The \emph{semantics} axis distinguishes how a heuristic condition reads partial information. Under the clingo-like semantics, conditions are evaluated with respect to \emph{established} truth values: a negative literal \mintinline{text}|not a| holds only if $a$ has been assigned false, and an aggregate refers to the value it has once its source atoms are completely determined. Under the Alpha-like semantics, conditions are evaluated \emph{optimistically} on the partial state: \mintinline{text}|not a| holds as soon as $a$ is not assigned true, and an aggregate is computed over the atoms currently true, so its value changes as search proceeds. The Alpha reading is what makes heuristics genuinely \emph{dynamic}: a directive such as \enquote{prefer the element whose insertion keeps the sums balanced} only makes sense if the sum is read off the evolving assignment.

Crossing the axes yields the four variants used throughout, summarized in Table~\ref{tab:variant-matrix}, plus the baseline \mintinline{text}|gc_noheur|, which contains the domain logic without any heuristic.

\begin{table}[H]
\centering
\begin{tabular}{lll}
\toprule
 & \textbf{clingo-like semantics} & \textbf{Alpha-like semantics}\\
\midrule
\textbf{Ground-and-solve} & \mintinline{text}|gc| (native \mintinline{text}|#heuristic|) & \mintinline{text}|ga| (simulated, see \S\ref{sec:ga-simulation})\\
\textbf{Lazy (propagator)} & \mintinline{text}|lc| & \mintinline{text}|la|\\
\bottomrule
\end{tabular}
\caption{The experimental matrix. Each cell exists for both evaluation backends (native and Prolog) and for each benchmark problem (BSP, PUP, HRP).}
\label{tab:variant-matrix}
\end{table}

Two structural constraints shape the matrix. First, ground clingo cannot express the Alpha semantics exactly: there is no ground construct denoting \enquote{the sum of the atoms currently true}. The \mintinline{text}|ga| variants therefore \emph{simulate} the Alpha reading, as described next, and only \mintinline{text}|la| captures it exactly. Second, the pair \mintinline{text}|gc|/\mintinline{text}|ga| is kept symmetric to the pair \mintinline{text}|lc|/\mintinline{text}|la|: within each approach, the two semantics differ exclusively in the treatment of negation and aggregates, so that any performance difference between suffixes can be attributed to the semantics and any difference between prefixes to the approach.

\subsubsection{Simulating Alpha in Ground clingo}
\label{sec:ga-simulation}

The \mintinline{text}|ga| encodings apply two systematic transformations to \mintinline{text}|gc|. The first transformation drops the negative literals on the partial state (for BSP, \mintinline{text}|not c(X)| and \mintinline{text}|not b(X)|; for PUP, the \mintinline{text}|not not_assigned_...| conditions; for HRP, the guards such as \mintinline{text}|not fullCabinet(C)|). In Alpha these literals would also be satisfied on undecided atoms, so removing them is the closest ground approximation of the Alpha reading; the solver anyway never re-decides atoms that are already assigned.

The second transformation replaces exact aggregate bindings by \emph{bounds over a static domain}. For BSP:

\begin{minted}{prolog}
sum_value(0..tot).
#heuristic b(X) : x(X), sum_value(S), S <= #sum{Y : c(Y)},
    W = ((n+1)*S)+X. [W, true]
\end{minted}

Instead of binding $S$ to the final value of the sum, the directive is unrolled over all candidate values \mintinline{text}|sum_value(S)| and guarded by the monotone condition $S \le \#\mathtt{sum}\{Y : c(Y)\}$. During search, as soon as the partial sum reaches a bound $S$, every directive with that bound becomes applicable; among the applicable ground directives for the same atom, clasp's Domain heuristic retains the one with the maximum weight, which corresponds exactly to the \emph{current} value of the growing aggregate. The max-weight selection thus reconstructs the dynamic aggregate at solving time. For aggregates whose contribution to the weight is decreasing rather than increasing (the free-capacity terms $N_0,N_1,N_2$ in PUP), the same trick is applied with the dual bound $V \ge \#\mathtt{max}\{\ldots\}$, so that the max-weight selection picks the smallest proven value. The price of the simulation is precisely the combinatorial explosion measured in Section~\ref{sec:results}: one ground directive per point of the product of the unrolled domains.

\subsubsection{Weights, Priorities, and Flattening}
\label{sec:flattening}

Alpha and clingo interpret the annotation $[W@P]$ differently, and the difference is central to this work. In Alpha, heuristic annotations $(W,P)$ form global levels: every applicable heuristic with priority $P$ strictly precedes every heuristic with priority $P' < P$, regardless of weights, and the weight discriminates within a level. In clingo, the priority \mintinline{text}|@P| only arbitrates between multiple heuristic modifications \emph{of the same atom}; the global ranking among different atoms is induced by the resolved level value, that is, by the weight alone.

Encodings written for Alpha therefore cannot be transplanted into clingo verbatim. Whenever an ordering between different atom families is intended (zones before sensors in PUP; reuse before cabinet-filling before room-filling in HRP), the levels must be \emph{flattened into the weight}:
\[
W_{\mathit{flat}} \;=\; W + P \cdot M,
\qquad M > \max W + 1,
\]
where $M$ is a level multiplier strictly larger than any intra-level weight, so that no heuristic of a lower level can ever outrank a heuristic of a higher level. In the PUP encodings, $M$ is the constant offset $100{,}000$ added to the zone weights (the maximum sensor weight is bounded by $46{,}220$ on the benchmark instances); in HRP, $M$ is computed from the instance as \mintinline{text}|levelMul(LM) :- LM = MC + MR + 10| over the maximum cabinet and room identifiers.

Flattening is order-preserving, and this is what makes it an admissible translation rather than an approximation. If $M$ exceeds every intra-level weight, then $W + P\cdot M > W' + P'\cdot M$ holds exactly when $(P,W)$ precedes $(P',W')$ lexicographically, so the total order induced on the atoms is the same one Alpha's absolute levels induce. Nothing of the strict discrimination between levels is lost; only its carrier changes, from a separate field to the weight itself.

The suite therefore adopts a single convention, applied uniformly to \emph{all four variants and both backends}: the absolute levels always live in the weight, and the priority field is reserved for its local role, that is, arbitrating between several directives attached to the \emph{same} atom, exactly as \mintinline{text}|@P| does in clingo. In BSP the priority is $0$ throughout, since there are no competing directives; in PUP likewise; in HRP the priority takes the two values $1$ and $0$, which suffice to let the \mintinline{text}|true| directives win over the closing \mintinline{text}|false| directives on the same target, the role that clingo's implicit default $P=|W|$ plays in the ground variants.

This convention is a deliberate methodological choice, not a limitation of the mechanism. The native propagator does implement Alpha's absolute levels natively: under \mintinline{text}|__semantics(alpha)| it ranks candidates globally by (priority, weight), and the C++ test suite covers exactly this behaviour, together with its local-only counterpart under \mintinline{text}|__semantics(clingo)|. The benchmark encodings do not exercise it, because the experiment requires the heuristic to be a controlled constant: if \mintinline{text}|la| expressed the order through native levels while \mintinline{text}|gc|, \mintinline{text}|ga|, and \mintinline{text}|lc| expressed it through flattened weights, a difference in results could not be attributed to the axes under study. Since ground clingo cannot rank different atoms by priority at all, flattening is the only mechanism available in every cell of the design, and hence the only one that keeps the four variants comparable. The same argument applies across backends: the Prolog backend sorts globally by (priority, weight) irrespective of the semantics selector, so an encoding that carried its levels in the priority field would not mean the same thing in the two backends. Under the convention, every encoding pair differs by exactly one token, \mintinline{text}|__semantics(alpha)| versus \mintinline{text}|__semantics(clingo)| in the native DSL and \mintinline{text}|alpha_not| versus \mintinline{text}|clingo_not| in the Prolog rules.

A final convention concerns empty aggregates: \mintinline{text}|#max| over an empty set is $0$ in Alpha but $\#\mathit{inf}$ in clingo. All encodings therefore use \mintinline{text}|#max{0; ...}| with an explicit neutral element, and both lazy backends implement the same convention (the native \mintinline{text}|MaxState| returns $0$ when empty; the Prolog runtime defines \mintinline{text}|dyn_max| with a $0$ default).

\subsection{System Architecture}

The project maintains two modified copies of clingo with the same source layout. The build in \path{clingo-native} implements the native backend; the build in \path{clingo-prolog} implements the Prolog backend and is configured with \mintinline{text}|CLINGO_USE_SWIPL=ON| so that the SWI-Prolog engine is linked in-process. In both builds the extension lives entirely under \path{libclingo} and is registered in \path{clingo_app.cc}:

\begin{minted}{cpp}
HeuristicPropagator my_propagator;
ctl.register_propagator(my_propagator, true);
\end{minted}

The class implements \mintinline{text}|Clingo::Heuristic|, the propagator interface that additionally exposes the \mintinline{text}|decide| callback. Both backends follow the same lifecycle: during \mintinline{text}|init| they scan the symbolic atoms for their trigger facts (\mintinline{text}|__heuristic(...)| in the native build, \mintinline{text}|heuristic("...")| in the Prolog build), build their runtime representation, and register watches on the solver literals they need to observe; during search, \mintinline{text}|propagate| and \mintinline{text}|undo| keep the internal state synchronized with the trail; at each decision point, \mintinline{text}|decide| either returns a signed literal for the best applicable heuristic target or returns $0$, leaving the decision to the solver's fallback heuristic. Runs are launched with \mintinline{text}|--heuristic=Domain|. Heuristics never change the answer sets: the propagator adds no constraints and only influences the order in which the solver explores the search space.

\subsection{The Native Backend}

\subsubsection{Declarative Syntax}

A native lazy heuristic is one ground fact per heuristic rule. The BSP heuristic, whose native directive would ground into $\Theta(n^3)$ objects, becomes:

\begin{minted}{prolog}
__heuristic(__target(b), x, __n_c, __bind(s, __sum(c)),
            __weight(__add(__mul(s, B), self)), __priority(0),
            __semantics(alpha), __modifier(true)) :- B = n + 1.
\end{minted}

The arguments are, in order: the target predicate; an implicit positive body predicate (\mintinline{text}|x|, matched on the target tuple); a negative body predicate (\mintinline{text}|__n_c| encodes \mintinline{text}|not c(X)| on the target tuple); an aggregate binding that introduces the runtime variable \mintinline{text}|s| as the incrementally maintained sum over the source predicate \mintinline{text}|c|; the weight expression, built from the operators \mintinline{text}|__add|, \mintinline{text}|__sub|, \mintinline{text}|__mul| over constants, bound variables, and the reserved symbol \mintinline{text}|self| (the first numeric argument of the target); the priority expression; the semantics selector; and the modifier. The ASP body \mintinline{text}|:- B = n + 1| is a small but important idiom: the fact is ground, and gringo evaluates the constant $B$ and injects it as a plain number, so the propagator never needs to know program constants.

Richer bodies are expressed explicitly. When a positive body predicate does not share the target tuple, \mintinline{text}|__body(pred, __match(src, tgt), __bind_arg(var, idx))| specifies the argument correspondence and extracts values into runtime variables; \mintinline{text}|__bind_target_arg(var, idx)| names a target argument. Aggregates support per-target filtering with offsets, as in

\begin{minted}{prolog}
__bind(m1, __max(num_sensors_on_unit, 0, __filter(1,1,-1)))
\end{minted}

which, for a target \mintinline{text}|assigned_zone_unit(Z,U)|, maintains the maximum of the first argument of \mintinline{text}|num_sensors_on_unit(N,U')| restricted to source atoms with $U' = U - 1$. Since the aggregate machinery draws from a single source predicate, joins occurring inside Alpha-style aggregates are precompiled as auxiliary derived predicates in the encoding (for PUP, \mintinline{text}|sensor_assigned_zone/2| and \mintinline{text}|sensor_zone_on_unit/3|); these auxiliaries appear in no constraint, so they cannot change the answer sets.

\subsubsection{Parser and Runtime Representation}

The parser in \path{libclingo/src/heuristic_parser.cc} turns each \mintinline{text}|__heuristic| symbol into a \mintinline{text}|HeuristicRuleTemplate|: target predicate, positive body specifications, negative body predicates, variable bindings, modifier, semantics, and two arithmetic ASTs for weight and priority. Malformed syntax, duplicate variables, unknown operators, or negative indices are rejected eagerly with specific error messages, so that a misspelled heuristic fails fast instead of silently guiding nothing.

During \mintinline{text}|init|, the propagator builds an index from predicate names to the ground literals of their numeric tuples, and \emph{materializes candidates}: for every ground target atom, and for every combination of matching positive body atoms, it creates a \mintinline{text}|RuntimeHeuristicCandidate| storing the target literal, the target tuple, the body literals, the values bound from body arguments, and the instantiated aggregate keys. Watches are registered on both polarities of the target and negative body literals, and on the positive body and aggregate source literals. Candidate materialization is proportional to the number of ground targets, but crucially it happens inside the propagator's memory as small POD structures, not as printed ground directives with their symbolic weights, and it does not multiply with the value domains of the aggregates.

\subsubsection{Incremental Aggregates}

Aggregate states implement a minimal interface with \mintinline{text}|add|, \mintinline{text}|remove|, and \mintinline{text}|result|. Sums and counts are plain accumulators; minima and maxima are ordered multisets, so that removal on backtracking is exact. Each source atom contributes its numeric value to the \mintinline{text}|RuntimeAggregateKey| obtained by instantiating the filter values; contributions are applied in \mintinline{text}|propagate| when a source literal becomes true and reverted in \mintinline{text}|undo|. Aggregates read under the clingo semantics are additionally \emph{gated}: the propagator tracks the full set of source literals per runtime key and considers the aggregate usable only when every source is determined, which is the clingo-like reading of aggregate values. Under the Alpha semantics the current value is always usable.

\subsubsection{Evaluation and Decision}

A candidate is active when its target is still free, all positive body literals are true, every negative body literal is satisfied according to the template's semantics,

\begin{minted}{cpp}
static bool negative_literal_is_satisfied(HeuristicSemantics s, Clingo::TruthValue v) {
    return s == HeuristicSemantics::Clingo ? v == Clingo::TruthValue::False
                                           : v != Clingo::TruthValue::True;
}
\end{minted}

and its variable environment can be built, which for clingo-semantics templates includes the aggregate gating above. Active effects are folded per target into a \mintinline{text}|TargetHeuristicState|. The modifiers follow clingo's Domain heuristic: \mintinline{text}|level| sets the ranking value, \mintinline{text}|sign| the preferred polarity, \mintinline{text}|true|/\mintinline{text}|false| set both; \mintinline{text}|init| and \mintinline{text}|factor| are recognized but rejected, because they modify solver-internal scores that the \mintinline{text}|decide| interface cannot reach, and mapping them to \mintinline{text}|level| would be semantically misleading.

The per-target resolution and the global ranking implement the two readings of priority discussed in Section~\ref{sec:flattening}. For clingo-semantics effects, the local priority arbitrates between modifications of the same target, and the target enters the global ranking with its resolved weight only. For Alpha-semantics effects, the best effect per target is selected by (priority, weight), and the pair itself becomes the global rank of the target, reproducing Alpha's levels. Ranked targets live in an ordered set:

\begin{minted}{cpp}
std::set<DecisionRankKey, DecisionRankGreater> active_decision_ranks_;
\end{minted}

ordered by decreasing priority, then decreasing weight, then a deterministic literal tie-break. \mintinline{text}|decide| pops stale entries lazily (targets that have been assigned, or whose state changed since ranking), applies the resolved sign to the best valid target, and otherwise falls back to the solver. \mintinline{text}|undo| uses \mintinline{text}|noexcept| refresh paths throughout, since exceptions during backtracking would violate the solver's expectations.

\subsection{The Prolog Backend}

\subsubsection{Rule Strings and Normalization}

In the Prolog backend, a heuristic is a rule string whose head is \mintinline{text}|heuristic(Target, Weight, Priority, Modifier)|:

\begin{minted}{prolog}
heuristic("heuristic(b(X), W, 0, true) :-
    aggregate_all(sum(Y), holds(c(Y)), S), holds(heuristic_base(Base)),
    holds(x(X)), target_available(b(X)), alpha_not(c(X)),
    W is Base*S+X.").
\end{minted}

The body is genuine Prolog. The propagator collects all \mintinline{text}|heuristic/1| facts (the alias \mintinline{text}|prolog_heuristic/1| is accepted for compatibility), normalizes each string, and classifies its semantics syntactically: a rule containing \mintinline{text}|clingo_not(...)| is clingo-like, otherwise it is Alpha-like. To keep the synchronized state small, the propagator statically extracts the predicate signatures that actually occur in the rules (inside \mintinline{text}|holds|, \mintinline{text}|alpha_not|, \mintinline{text}|clingo_not|, \mintinline{text}|target_available|, and as rule targets) and synchronizes only atoms with those signatures.

\subsubsection{The Runtime Program}

At initialization the backend boots the SWI-Prolog engine in-process \cite{wielemaker2012swipl} and consults a generated runtime program that defines the evaluation vocabulary:

\begin{minted}{prolog}
holds(A) :- true_atom(A).
holds(A) :- static_atom(A), \+ true_atom(A).
clingo_not(A) :- false_atom(A).
alpha_not(A)  :- \+ holds(A).
target_available(A) :- target_atom(A), \+ true_atom(A), \+ false_atom(A).
dyn_max(Goal, Template, Max) :-
    (aggregate_all(max(Template), Goal, R) -> Max = R ; Max = 0).
\end{minted}

These few clauses are the entire semantic bridge. The solver trail is mirrored by the dynamic predicates \mintinline{text}|true_atom/1| and \mintinline{text}|false_atom/1|; a free atom is represented by absence from both. \mintinline{text}|alpha_not| and \mintinline{text}|clingo_not| implement exactly the two negation readings of Section~\ref{sec:ga-simulation}; \mintinline{text}|aggregate_all| over \mintinline{text}|holds| computes dynamic aggregates on the atoms currently true, which is the Alpha aggregate semantics; \mintinline{text}|dyn_max|/\mintinline{text}|dyn_min| add the empty-set-is-zero convention; \mintinline{text}|target_available| guarantees that the rule only proposes targets that are still free.

\subsubsection{Synchronization and Decision}

Trail synchronization is incremental: \mintinline{text}|propagate| and \mintinline{text}|undo| update only the atoms mapped by the changed literals, each with a single combined retract-and-assert goal, since term parsing and calling dominate the synchronization cost. At each \mintinline{text}|decide|, the backend enumerates the solutions of \mintinline{text}|heuristic(T, W, P, M)|, discards candidates whose target is assigned or unmapped, and selects the best candidate by decreasing priority, then decreasing weight, then a deterministic tie-break; the modifier gives the sign of the returned literal. The backend is enabled by the environment variable \mintinline{text}|LAZY_HEURISTIC_BACKEND=prolog|; with \mintinline{text}|LAZY_PROLOG_STATS=1| it prints a summary line with decision counts and cumulative timings of every phase (state synchronization, Prolog query, candidate scan, selection), which the benchmark parser ingests.

Without the environment variable, the same build falls back to an auxiliary evaluator that re-grounds a small clingo program per decision. This fallback accepts a different, ASP-style body syntax; a Prolog-style rule evaluated by it simply fails (SWI's \mintinline{text}|unknown=fail| convention), so the run silently degenerates to no heuristic at all. This failure mode is invisible in exit codes and therefore dangerous for benchmarking; the guards that detect it are described in Section~\ref{sec:validation}.

\subsection{Benchmark Problems and Encodings}

\subsubsection{Balanced Set Partitioning (BSP)}

BSP partitions $\{1,\ldots,n\}$ into two sets \mintinline{text}|b/1| and \mintinline{text}|c/1| whose sums differ by at most one. The guessing part is a symmetric even loop, the check compares the two sums, and the heuristic prefers to place the next element into the set whose current sum is largest, thereby keeping the partition balanced greedily. BSP is also the one benchmark on which the flattening rule of Section~\ref{sec:flattening} is vacuous: all its heuristic directives belong to a single level, since there is no ordering between distinct atom families to encode (\mintinline{text}|b/1| and \mintinline{text}|c/1| are ranked by one and the same criterion). The weight $((n+1)\cdot S)+X$ is therefore not a flattened weight--priority pair but the lexicographic composition of the two components of a single criterion, the dynamic sum $S$ dominating and the element index $X$ breaking ties inside it; the multiplier $n+1$ plays the role of the level multiplier $M$ within the criterion rather than between levels. Accordingly all four variants carry priority $0$, including \mintinline{text}|la|, and the \mintinline{text}|la|/\mintinline{text}|lc| contrast on BSP isolates the semantics axis alone, with no interference from the priority reading. Besides the four matrix variants and the baseline, BSP has three methodological extras: \mintinline{text}|ga_weak|, which drops the negative literals without the static unrolling and therefore changes only the activation moment; \mintinline{text}|la_aux|, which routes the lazy heuristic through an auxiliary predicate \mintinline{text}|h_b(X,S)| that reintroduces grounded aggregate values (and hence a grounding cost) before the propagator is even used; and \mintinline{text}|la_co|, which combines the lazy heuristic with a linear reformulation of the balance constraint, and is treated as a modelling experiment rather than a heuristic comparison.

\subsubsection{Partner Units Problem (PUP)}

PUP \cite{aschinger2011pup} connects zones and sensors to communication units under adjacency and capacity constraints; the instances (the \mintinline{text}|double| family; the related \mintinline{text}|doubleV| family is left to future work) provide \mintinline{text}|comUnit/1| and \mintinline{text}|zone2sensor/2|. The encoding follows the Alpha evaluation encoding with an inductive, locality-based choice structure. The heuristic ranks sensor placements by a four-component weight (dynamic degree, forbidden placements, unit occupancy, direct connections) and zone placements by occupancy and free capacity, with zones globally preceding sensors.

Porting this encoding from Alpha to clingo required documented adaptations, applied uniformly to \emph{all} variants so that comparisons remain fair: the inductive \mintinline{text}|assignable| rules are bounded by \mintinline{text}|comUnit(U±1)|, because gringo would otherwise materialize an unbounded chain that Alpha explores lazily; the existential capacity constraints are rewritten as monotone counting aggregates with linear grounding; a redundant distance-two blocking rule with cubic grounding is removed; counters are capped at the unit capacity (\mintinline{text}|sensorNumbers(0..2)|), without which the auxiliary predicate \mintinline{text}|num_forbidden_places_of_sensors| would explode in every variant and mask the effect under study. The zones-before-sensors ordering is flattened into the weight in all four variants and in both backends (offset $+100{,}000$, against a maximum sensor weight of $46{,}220$), with the priority field left at $0$ throughout, per the convention of Section~\ref{sec:flattening}.

\subsubsection{House Reconfiguration Problem (HRP)}

HRP \cite{friedrich2011reconfiguration} extends the House Configuration Problem with legacy configurations: things must be stored in cabinets and cabinets placed in rooms owned by the things' owners, under capacity constraints refined by the thing-length/cabinet-size extension, while a legacy configuration can be partially reused or dismantled. The heuristic is a strict five-level policy: reuse legacy components first (with internal weights), then fill cabinets with long things, then with short things, then place cabinets into rooms, and finally a closing level that biases all remaining choice atoms to false, pruning the tail of the search. HRP is the negation-heavy benchmark: its heuristics contain no aggregates, so the \mintinline{text}|la|/\mintinline{text}|lc| comparison isolates the effect of the negation semantics alone, while in BSP and PUP both mechanisms act together. In the native lazy encodings, the guards over the partial state are precompiled into same-tuple auxiliary predicates (for example \mintinline{text}|cabFull(C,T) :- cabinetDomain(C), thing(T), fullCabinet(C)|) so that negative conditions align with the target tuple; the five levels are flattened with the instance-derived multiplier \mintinline{text}|levelMul(LM) :- LM = MC + MR + 10| in all four variants and in both backends, and the priority field retains only its local role, distinguishing the \mintinline{text}|true| directives from the closing \mintinline{text}|false| ones on the same target.

\subsection{Correctness Validation}
\label{sec:validation}

Heuristics must not change the answer sets: this is a falsifiable property, and the pipeline verifies it before measuring anything. The harness \path{tools/check_equivalence.py}, run by the sanity entry point ahead of every campaign, performs an \emph{exhaustive equivalence} check on a small BSP instance: it enumerates all models with \mintinline{text}|clingo -n 0| for every variant and both backends, projects the models onto the solution predicates (\mintinline{text}|b/1|, \mintinline{text}|c/1|), and requires the projected collections to be identical across variants and across backends. Projection is necessary because the encodings expose different \mintinline{text}|#show| sets; comparing raw witnesses would produce false positives. For PUP, where exhaustive enumeration is intractable even on small instances, the harness checks agreement on SAT/UNSAT and reports the choice counts.

The harness also guards against the silent-fallback failure mode of the Prolog backend. Every lazy Prolog run must report \mintinline{text}|decide_calls > 0| in the propagator summary: a run that terminated successfully but never called the custom \mintinline{text}|decide| has degenerated to the no-heuristic baseline and is treated as an error by the harness (and flagged with an explicit warning column by the benchmark runner). As a cross-check, the harness verifies that native \mintinline{text}|la| and Prolog \mintinline{text}|la| perform the same number of choices on the control instance: if the SWI backend were inactive, the Prolog run would diverge from the native one. The benchmark wrapper for the Prolog system exports \mintinline{text}|LAZY_HEURISTIC_BACKEND=prolog| itself, so the guard protects against configuration drift rather than repairing it. The C++ unit tests complete the picture on the native side: they cover incremental aggregates, corner cases (empty domains, unsatisfiable instances), explicit body and target bindings, filtered aggregates, the local-versus-global reading of priorities in the two semantics, and the syntax validation of the parser.

\subsection{Benchmark Infrastructure and Metrics}

The campaigns are orchestrated with \mintinline{text}|benchmark-tool|, driving one run per (system, setting, instance) triple under \mintinline{text}|runlim| \cite{runlim} for wall-clock and memory limits. The two systems are shell wrappers around the two modified binaries; both fix \mintinline{text}|--stats=2 --heuristic=Domain -n 1| and a numeric locale, and the Prolog wrapper additionally exports the backend activation and statistics variables. A custom result parser extracts, for every run: outcome and times (total, solving, and their difference as an estimate of the non-solving part), search statistics (choices, conflicts, restarts), grounding-size measures (ground \mintinline{text}|#heuristic| directives for the ground variants, lazy heuristic facts for the lazy ones, and rule counts from clasp's statistics), the peak memory, and, for the Prolog backend, the propagator summary (decision calls, candidates seen per decision, cumulative synchronization, query, scan, and selection times).

Two measurement caveats are declared up front, because they define what the numbers can and cannot show. First, the memory metric is the peak resident set size of the whole solver process, end to end: for the lazy Prolog variants it includes the in-process SWI engine and its asserted database. The engine has a fixed initialization overhead independent of $n$, so on small instances the lazy variants may show a \emph{larger} RSS than the ground ones; the lazy memory advantage is asymptotic, emerging where the $\Theta(n^3)$ growth of ground heuristics overtakes the constant overhead. Second, the ground-size comparison counts objects of different kinds: for ground variants it counts directives actually materialized by gringo (a faithful one-line-one-object count), whereas for lazy variants it counts template facts that expand at runtime into one candidate per ground target. The line count therefore \emph{understates} the runtime work of the lazy backends by design; that work is measured separately by the per-decision metrics. Both caveats are taken up again in the discussion of the results.
